% GENERATED FILE. Edit canonical lessons, then run npm run generate:books.
% canonical-source-hash: fnv1a64:fab3c7c98e2cf3a9
% canonical-lessons: KA-C429-asisu, KA-C429-apeksisu, KA-C429-prastapisu, KA-C429-sucisu, KA-C429-parihara

\chapter{Hopes and Proposals}
\label{ch:ka-hopes-and-proposals}
\hlchaptermodality{429}

\begin{chapteropening}
\textbf{By the end of this chapter:} \emph{I can say to hope for, to look to, to propose, to suggest and a solution.}

Complete the last lesson of chapter 429: I can say to hope for, to look to, to propose, to suggest and a solution.
\end{chapteropening}

\section[\texorpdfstring{\kn{ಆಶಿಸು} (āśisu)}{ಆಶಿಸು (āśisu)}]{\kn{ಆಶಿಸು} (āśisu) — to hope for}

\label{lesson:KA-C429-asisu}



\begin{quote}
Before the new one: say the Kannada for applicability, then the Kannada for overall.
\end{quote}

\subsection*{You'll want to know: \kn{ಆಶಿಸು}}
\textbf{\kn{ಆಶಿಸು}} — \emph{āśisu} — \textquotedblleft{}to hope for\textquotedblright{}.

\begin{quote}
\kn{ಮಗ} \kn{ಪರೀಕ್ಷೆಯಲ್ಲಿ} \kn{ಪಾಸಾಗುತ್ತಾನೆ} \kn{ಅಂತ} \kn{ತಾಯಿ} \kn{ಆಶಿಸಿದಳು}.
\end{quote}

\begin{quote}
\emph{maga parīkṣeyalli pāsāguttāne anta tāyi āśisidaḷu.}
\end{quote}

\begin{quote}
The mother hoped her son would pass the exam.
\end{quote}

Past (I …): \kn{ಆಶಿಸಿದೆ}, \emph{āśiside}.

What is hoped comes before \kn{ಅಂತ} or \kn{ಎಂದು}.

\begin{cousinweb}
\textbf{\kn{ಆಶಿಸು}} (\emph{āśisu}) is built on Sanskrit \emph{āśā}, \textquotedblleft{}wish, desire, hope\textquotedblright{} (Monier-Williams), with the Kannada verb ending -\emph{isu}.
\end{cousinweb}

\begin{grammarlens}[title={What you've built}]
One more word for saying what happens if.
\end{grammarlens}

\subsection*{Your turn}
\begin{itemize}
\raggedright
  \item \emph{Say it:} \emph{āśisu}
  \item \emph{Say it:} \emph{āśisu}, once more
  \item \emph{Say it:} \emph{maga parīkṣeyalli pāsāguttāne anta tāyi āśisidaḷu.} — The mother hoped her son would pass the exam.
  \item \emph{Recall:} say the Kannada for applicability, then the Kannada for overall, then say \emph{āśisu} again
\end{itemize}

\begin{tcolorbox}[breakable,colback=teal!4,colframe=teal!35!black,title={Before you move on}]
What does \kn{ಆಶಿಸು} mean? (\textquotedblleft{}To hope for\textquotedblright{}.) Say it once more.
\end{tcolorbox}
\section[\texorpdfstring{\kn{ಅಪೇಕ್ಷಿಸು} (apēkṣisu)}{ಅಪೇಕ್ಷಿಸು (apēkṣisu)}]{\kn{ಅಪೇಕ್ಷಿಸು} (apēkṣisu) — to look to; to expect}

\label{lesson:KA-C429-apeksisu}



\begin{quote}
Before the new one: say the Kannada for overall, then the Kannada for to hope for.
\end{quote}

\subsection*{You'll want to know: \kn{ಅಪೇಕ್ಷಿಸು}}
\textbf{\kn{ಅಪೇಕ್ಷಿಸು}} — \emph{apēkṣisu} — \textquotedblleft{}to look to; to expect\textquotedblright{}.

\begin{quote}
\kn{ಅವನು} \kn{ಯಾರಿಂದಲೂ} \kn{ಸಹಾಯ} \kn{ಅಪೇಕ್ಷಿಸಲಿಲ್ಲ}.
\end{quote}

\begin{quote}
\emph{avanu yārindalū sahāya apēkṣisalilla.}
\end{quote}

\begin{quote}
He expected help from no one.
\end{quote}

Past (I …): \kn{ಅಪೇಕ್ಷಿಸಿದೆ}, \emph{apēkṣiside}.

\kn{ಅಪೇಕ್ಷೆ}, a wish or an expectation.

\begin{cousinweb}
\textbf{\kn{ಅಪೇಕ್ಷಿಸು}} (\emph{apēkṣisu}) is built on Sanskrit \emph{apekṣā}, \textquotedblleft{}looking for, expectation, need\textquotedblright{} (Monier-Williams), with the Kannada verb ending -\emph{isu}.
\end{cousinweb}

\begin{grammarlens}[title={What you've built}]
One more word for saying what happens if.
\end{grammarlens}

\subsection*{Your turn}
\begin{itemize}
\raggedright
  \item \emph{Say it:} \emph{apēkṣisu}
  \item \emph{Say it:} \emph{apēkṣisu}, once more
  \item \emph{Say it:} \emph{avanu yārindalū sahāya apēkṣisalilla.} — He expected help from no one.
  \item \emph{Recall:} say the Kannada for overall, then the Kannada for to hope for, then say \emph{apēkṣisu} again
\end{itemize}

\begin{tcolorbox}[breakable,colback=teal!4,colframe=teal!35!black,title={Before you move on}]
What does \kn{ಅಪೇಕ್ಷಿಸು} mean? (\textquotedblleft{}To look to; to expect\textquotedblright{}.) Say it once more.
\end{tcolorbox}
\section[\texorpdfstring{\kn{ಪ್ರಸ್ತಾಪಿಸು} (prastāpisu)}{ಪ್ರಸ್ತಾಪಿಸು (prastāpisu)}]{\kn{ಪ್ರಸ್ತಾಪಿಸು} (prastāpisu) — to propose; to bring up}

\label{lesson:KA-C429-prastapisu}



\begin{quote}
Before the new one: say the Kannada for to hope for, then the Kannada for to look to.
\end{quote}

\subsection*{You'll want to know: \kn{ಪ್ರಸ್ತಾಪಿಸು}}
\textbf{\kn{ಪ್ರಸ್ತಾಪಿಸು}} — \emph{prastāpisu} — \textquotedblleft{}to propose; to bring up\textquotedblright{}.

\begin{quote}
\kn{ಸಭೆಯಲ್ಲಿ} \kn{ಅವಳು} \kn{ಹೊಸ} \kn{ಯೋಜನೆಯನ್ನು} \kn{ಪ್ರಸ್ತಾಪಿಸಿದಳು}.
\end{quote}

\begin{quote}
\emph{sabheyalli avaḷu hosa yōjaneyannu prastāpisidaḷu.}
\end{quote}

\begin{quote}
At the meeting she proposed a new scheme.
\end{quote}

Past (I …): \kn{ಪ್ರಸ್ತಾಪಿಸಿದೆ}, \emph{prastāpiside}.

\kn{ಪ್ರಸ್ತಾಪ}, a proposal.

\begin{cousinweb}
\textbf{\kn{ಪ್ರಸ್ತಾಪಿಸು}} (\emph{prastāpisu}) is built on Sanskrit \emph{prastāva}, \textquotedblleft{}an introduction, a prelude\textquotedblright{} (Monier-Williams), with the Kannada verb ending -\emph{isu}.
\end{cousinweb}

\begin{grammarlens}[title={What you've built}]
One more word for saying what happens if.
\end{grammarlens}

\subsection*{Your turn}
\begin{itemize}
\raggedright
  \item \emph{Say it:} \emph{prastāpisu}
  \item \emph{Say it:} \emph{prastāpisu}, once more
  \item \emph{Say it:} \emph{sabheyalli avaḷu hosa yōjaneyannu prastāpisidaḷu.} — At the meeting she proposed a new scheme.
  \item \emph{Recall:} say the Kannada for to hope for, then the Kannada for to look to, then say \emph{prastāpisu} again
\end{itemize}

\begin{tcolorbox}[breakable,colback=teal!4,colframe=teal!35!black,title={Before you move on}]
What does \kn{ಪ್ರಸ್ತಾಪಿಸು} mean? (\textquotedblleft{}To propose; to bring up\textquotedblright{}.) Say it once more.
\end{tcolorbox}
\section[\texorpdfstring{\kn{ಸೂಚಿಸು} (sūcisu)}{ಸೂಚಿಸು (sūcisu)}]{\kn{ಸೂಚಿಸು} (sūcisu) — to suggest; to indicate}

\label{lesson:KA-C429-sucisu}



\begin{quote}
Before the new one: say the Kannada for to look to, then the Kannada for to propose.
\end{quote}

\subsection*{You'll want to know: \kn{ಸೂಚಿಸು}}
\textbf{\kn{ಸೂಚಿಸು}} — \emph{sūcisu} — \textquotedblleft{}to suggest; to indicate\textquotedblright{}.

\begin{quote}
\kn{ವೈದ್ಯರು} \kn{ಸ್ವಲ್ಪ} \kn{ವಿಶ್ರಾಂತಿ} \kn{ಸೂಚಿಸಿದರು}.
\end{quote}

\begin{quote}
\emph{vaidyaru svalpa viśrānti sūcisidaru.}
\end{quote}

\begin{quote}
The doctor suggested some rest.
\end{quote}

Past (I …): \kn{ಸೂಚಿಸಿದೆ}, \emph{sūciside}.

\kn{ಸೂಚನೆ}, a notice, is the noun beside it.

\begin{cousinweb}
\textbf{\kn{ಸೂಚಿಸು}} (\emph{sūcisu}) is built on Sanskrit \emph{sūcana}, \textquotedblleft{}pointing out, indicating\textquotedblright{} (Monier-Williams), with the Kannada verb ending -\emph{isu}.
\end{cousinweb}

\begin{grammarlens}[title={What you've built}]
One more word for saying what happens if.
\end{grammarlens}

\subsection*{Your turn}
\begin{itemize}
\raggedright
  \item \emph{Say it:} \emph{sūcisu}
  \item \emph{Say it:} \emph{sūcisu}, once more
  \item \emph{Say it:} \emph{vaidyaru svalpa viśrānti sūcisidaru.} — The doctor suggested some rest.
  \item \emph{Recall:} say the Kannada for to look to, then the Kannada for to propose, then say \emph{sūcisu} again
\end{itemize}

\begin{tcolorbox}[breakable,colback=teal!4,colframe=teal!35!black,title={Before you move on}]
What does \kn{ಸೂಚಿಸು} mean? (\textquotedblleft{}To suggest; to indicate\textquotedblright{}.) Say it once more.
\end{tcolorbox}
\section[\texorpdfstring{\kn{ಪರಿಹಾರ} (parihāra)}{ಪರಿಹಾರ (parihāra)}]{\kn{ಪರಿಹಾರ} (parihāra) — a solution; compensation}

\label{lesson:KA-C429-parihara}



\begin{quote}
Before the new one: say the Kannada for to propose, then the Kannada for to suggest.
\end{quote}

\subsection*{You'll want to know: \kn{ಪರಿಹಾರ}}
\textbf{\kn{ಪರಿಹಾರ}} — \emph{parihāra} — \textquotedblleft{}a solution; compensation\textquotedblright{}.

\begin{quote}
\kn{ಈ} \kn{ಸಮಸ್ಯೆಗೆ} \kn{ಸುಲಭ} \kn{ಪರಿಹಾರ} \kn{ಇಲ್ಲ}.
\end{quote}

\begin{quote}
\emph{ī samasyege sulabha parihāra illa.}
\end{quote}

\begin{quote}
There is no easy solution to this problem.
\end{quote}

Also compensation: \kn{ಬೆಳೆ} \kn{ನಷ್ಟಕ್ಕೆ} \kn{ಪರಿಹಾರ}.

\begin{cousinweb}
\textbf{\kn{ಪರಿಹಾರ}} (\emph{parihāra}) is borrowed from Sanskrit \emph{parihāra}, \textquotedblleft{}leading round, avoiding; a remedy\textquotedblright{} (CDIAL 7901).
\end{cousinweb}

\begin{grammarlens}[title={What you've built}]
One more word for saying what happens if.
\end{grammarlens}

\subsection*{Your turn}
\begin{itemize}
\raggedright
  \item \emph{Say it:} \emph{parihāra}
  \item \emph{Say it:} \emph{parihāra}, once more
  \item \emph{Say it:} \emph{ī samasyege sulabha parihāra illa.} — There is no easy solution to this problem.
  \item \emph{Recall:} say the Kannada for to propose, then the Kannada for to suggest, then say \emph{parihāra} again
\end{itemize}

\begin{tcolorbox}[breakable,colback=teal!4,colframe=teal!35!black,title={Before you move on}]
What does \kn{ಪರಿಹಾರ} mean? (\textquotedblleft{}A solution; compensation\textquotedblright{}.) Say it once more.
\end{tcolorbox}
